\problem{4.46}{Three-dimensional harmonic oscillator and degeneracies}
For $V(\mathbf r)=\frac12m\omega^2(x^2+y^2+z^2)$,
the Hamiltonian is the sum of three commuting
one-dimensional oscillator Hamiltonians.
Product eigenstates
$\psi_{n_x}(x)\psi_{n_y}(y)\psi_{n_z}(z)$
therefore have energies
\begin{equation}
E=\sh\omega(n_x+n_y+n_z+\tfrac32),
\qquad n_x,n_y,n_z=0,1,\dots.
\end{equation}
Writing $N=n_x+n_y+n_z$,
\begin{equation}
\boxed{E_N=(N+\tfrac32)\sh\omega}.
\end{equation}
The number of ordered nonnegative triples
whose entries sum to $N$ is
$\binom{N+2}{2}$ (stars and bars). Hence
\begin{equation}
\boxed{d(N)=\frac{(N+1)(N+2)}2}.
\end{equation}
The first degeneracies are $1,3,6,10,\dots$.

\problem{4.47}{Spherical-coordinate solution of the isotropic oscillator}
Separate $\psi=R(r)Y_\ell^m(\theta,\phi)$,
and introduce
$x=\sqrt{m\omega/\sh}\,r$ and
$\epsilon=2E/(\sh\omega)$.
The radial equation is
\begin{equation}
R''+\frac2xR'+
\left[\epsilon-x^2-\frac{\ell(\ell+1)}{x^2}\right]R=0.
\end{equation}
Correct near-origin and asymptotic behavior
suggests $R=x^\ell\ee^{-x^2/2}v(x)$.
After differentiation and cancellation,
\begin{equation}
v''+\left(\frac{2\ell+2}{x}-2x\right)v'
+[\epsilon-(2\ell+3)]v=0.
\end{equation}
Put $v(x)=\sum_{j=0}^\infty c_jx^j$.
The recurrence is
\begin{equation}
\boxed{c_{j+2}=
\frac{2j+2\ell+3-\epsilon}
{(j+2)(j+2\ell+3)}\,c_j}.
\end{equation}
Regularity selects an even series.
Normalizability requires truncation at
$j=2n_r$, which gives
$\epsilon=4n_r+2\ell+3$ and
\begin{equation}
\boxed{E_{n_r\ell}
=(2n_r+\ell+\tfrac32)\sh\omega},
\qquad n_r,\ell=0,1,\dots.
\end{equation}
Identify $N=2n_r+\ell$ from Problem 4.46.
For a given $N$, $\ell=N,N-2,\dots\ge0$.
Summing the $2\ell+1$ magnetic states yields
\begin{equation}
\boxed{\sum_{\ell=N,N-2,\dots}(2\ell+1)
=\frac{(N+1)(N+2)}2}.
\end{equation}
This checks complete agreement of
Cartesian and spherical separations.

\problem{4.48}{Three-dimensional virial theorem}
Take $G=\mathbf r\cdot\mathbf p$ and use
$\dd\avg G/\dd t=(\ii/\sh)\avg{[H,G]}$.
The canonical commutators give
\begin{equation}
\boxed{\frac{\dd\avg G}{\dd t}
=2\avg T-\avg{\mathbf r\cdot\nabla V}}.
\end{equation}
For a stationary state $\dd\avg G/\dd t=0$
and therefore
\begin{equation}
\boxed{2\avg T=\avg{\mathbf r\cdot\nabla V}}.
\end{equation}
For Coulomb $V\propto r^{-1}$,
$\mathbf r\cdot\nabla V=-V$,
hence $2\avg T=-\avg V$.
Using $E_n=\avg T+\avg V$ yields
\begin{equation}
\boxed{\avg T=-E_n,\qquad\avg V=2E_n}.
\end{equation}
For $V\propto r^2$,
$\mathbf r\cdot\nabla V=2V$, giving
$\boxed{\avg T=\avg V=E_n/2}$.
These identities assume normalizable
states and vanishing surface terms.

\problem{4.49}{Probability current and hydrogen orbital flow}
The three-dimensional current is
\begin{equation}
\boxed{\mathbf J=\frac{\sh}{2m\ii}
(\Psi^*\nabla\Psi-\Psi\nabla\Psi^*)}.
\end{equation}
Take the Schrödinger equation
$\ii\sh\partial_t\Psi=
-\sh^2\nabla^2\Psi/(2m)+V\Psi$,
multiply by $\Psi^*$, and subtract the
complex-conjugate equation multiplied by
$\Psi$. The real-potential terms cancel:
\begin{equation}
\partial_t|\Psi|^2=
-\nabla\cdot\frac{\sh}{2m\ii}
(\Psi^*\nabla\Psi-\Psi\nabla\Psi^*)
=\boxed{-\nabla\cdot\mathbf J}.
\end{equation}
Integrating over a fixed volume gives
$\oint\mathbf J\cdot\dd\mathbf a
=-\dd(\int_V|\Psi|^2\dd^3r)/\dd t$.
For $\psi_{211}=-r\sin\theta
\ee^{-r/(2\aB)}\ee^{\ii\phi}/
(8\sqrt\pi\aB^{5/2})$, only
the azimuthal phase contributes:
\begin{equation}
\mathbf J=\frac{\sh}{mr\sin\theta}
|\psi_{211}|^2\hat{\boldsymbol\phi}
=\boxed{\frac{\sh}{64\pi m\aB^5}
r\ee^{-r/\aB}\sin\theta\,
\hat{\boldsymbol\phi}}.
\end{equation}
The mass-flow orbital angular momentum is
$m\int(\mathbf r\times\mathbf J)\dd^3r$.
Its $z$ component is
$m\int r\sin\theta\,J_\phi\dd^3r=
\sh\int|\psi_{211}|^2\dd^3r$, so
\begin{equation}
\boxed{L_z=\sh},
\end{equation}
as expected for $m_\ell=+1$.

\problem{4.50}{Momentum-space hydrogen ground state}
For
$\psi_{100}(\mathbf r)=
\ee^{-r/\aB}/\sqrt{\pi\aB^3}$,
define the normalized Fourier transform
\begin{equation}
\phi(\mathbf p)=\frac{1}{(2\pi\sh)^{3/2}}
\int_{\mathbb R^3}\ee^{-\ii\mathbf p\cdot
\mathbf r/\sh}\psi_{100}(\mathbf r)\dd^3r.
\end{equation}
Choose the polar axis along $\mathbf p$.
The angular integral is
$\int\dd\Omega\,\ee^{-\ii pr\cos\theta/\sh}
=4\pi\sin(pr/\sh)/(pr/\sh)$.
Using
$\int_0^\infty r\ee^{-\beta r}\sin(kr)\dd r
=2\beta k/(\beta^2+k^2)^2$, gives
\begin{equation}
\boxed{\phi(\mathbf p)=\frac1\pi
\left(\frac{2\aB}{\sh}\right)^{3/2}
\frac1{[1+(\aB p/\sh)^2]^2}}.
\end{equation}
Put $u=\aB p/\sh$. Its norm is
\begin{equation}
\int|\phi|^2\dd^3p
=\frac{32}{\pi}\int_0^\infty
\frac{u^2\dd u}{(1+u^2)^4}=1,
\end{equation}
where the integral equals $\pi/32$.
Similarly,
\begin{equation}
\avg{p^2}=\frac{32\sh^2}{\pi\aB^2}
\int_0^\infty\frac{u^4\dd u}{(1+u^2)^4}
=\boxed{\frac{\sh^2}{\aB^2}}.
\end{equation}
Thus $\boxed{\avg T=\sh^2/(2m\aB^2)=-E_1}$,
in agreement with the Coulomb virial theorem.

\problem{4.51}{Why every attractive circular well binds in two dimensions}
For a radially symmetric two-dimensional
potential $V=-V_0$ inside $r<a$
and zero outside, a bound $s$-wave has
$-V_0<E<0$, with
$k^2=2m(V_0+E)/\sh^2$ and
$\kappa^2=-2mE/\sh^2$.
The finite-at-origin interior solution is
$R_{\rm in}=AJ_0(kr)$; the decaying exterior
solution is $R_{\rm out}=BK_0(\kappa r)$.
Matching $R$ and $R'$ at $r=a$ gives
\begin{equation}
\boxed{k\frac{J_1(ka)}{J_0(ka)}
=\kappa\frac{K_1(\kappa a)}{K_0(\kappa a)}}.
\end{equation}
For arbitrarily weak attraction, $ka\ll1$.
As $\kappa a\to0$,
$K_0(\kappa a)\sim-\ln(\kappa a/2)-\gamma_E$,
$\kappa K_1(\kappa a)\sim1/a$, while
$kJ_1(ka)/J_0(ka)\sim k^2a/2$.
Hence
\begin{equation}
-\ln(\kappa a/2)-\gamma_E
\sim\frac{2}{(ka)^2},
\quad
\boxed{\kappa a\sim
2\exp[-\gamma_E-2/(ka)^2]}.
\end{equation}
For every $V_0>0$ this is nonzero,
though exponentially small, so there
is at least one negative-energy bound state.
Unlike three dimensions, there is no
nonzero threshold depth for an attractive
two-dimensional circular well.

\problem{4.52}{Hydrogen $\psi_{321}$ and radial moments}
For $n=3,\ell=2$,
$R_{32}=(2/(3\aB))^{3/2}\rho^2
\ee^{-\rho/2}/\sqrt{720}$,
$\rho=2r/(3\aB)$.
With
$Y_2^1=-\sqrt{15/(8\pi)}
\sin\theta\cos\theta\ee^{\ii\phi}$,
simplification gives
\begin{equation}
\boxed{\psi_{321}(r,\theta,\phi)=
-\frac{r^2}{81\sqrt\pi\,\aB^{7/2}}
\ee^{-r/(3\aB)}
\sin\theta\cos\theta\ee^{\ii\phi}}.
\end{equation}
For normalization, the angular integral is
$\int\sin^2\theta\cos^2\theta\dd\Omega=8\pi/15$;
the radial integral is
\begin{equation}
\int_0^\infty r^6\ee^{-2r/(3\aB)}\dd r
=6!\left(\frac{3\aB}{2}\right)^7.
\end{equation}
Combining with the displayed prefactor yields one.
For arbitrary real $s$,
\begin{equation}
\boxed{\avg{r^s}=
\frac{\Gamma(s+7)}{\Gamma(7)}
\left(\frac{3\aB}{2}\right)^s},
\qquad\boxed{s>-7}.
\end{equation}
The restriction follows from the
$r^{s+6}$ behavior near $r=0$.

\problem{4.53}{Hydrogen $\psi_{433}$ and transverse angular momentum}
For $n=4,\ell=3,m=3$,
$R_{43}=(1/(2\aB))^{3/2}\rho^3
\ee^{-\rho/2}/\sqrt{8\,7!}$,
where $\rho=r/(2\aB)$, and
$Y_3^3=-\sqrt{35/(64\pi)}
\sin^3\theta\ee^{3\ii\phi}$.
Therefore
\begin{equation}
\boxed{\psi_{433}=
-\frac{r^3}{6144\sqrt\pi\,\aB^{9/2}}
\ee^{-r/(4\aB)}
\sin^3\theta\ee^{3\ii\phi}}.
\end{equation}
The exact hydrogen radial expectation is
$\avg r=\frac{\aB}{2}[3n^2-\ell(\ell+1)]$;
here
\begin{equation}
\boxed{\avg r=18\aB}.
\end{equation}
Since $L_x^2+L_y^2=L^2-L_z^2$,
the state is already an eigenstate of the
requested observable:
\begin{equation}
\boxed{(L_x^2+L_y^2)\psi_{433}
=\sh^2[3(4)-3^2]\psi_{433}
=3\sh^2\psi_{433}}.
\end{equation}
A measurement returns $3\sh^2$ with probability one.

\problem{4.54}{Ground-state penetration of a finite-size nucleus}
For the hydrogen $1s$ state,
the probability inside a sphere of radius $b$ is
\begin{align}
P(r<b)&=\frac4{\aB^3}\int_0^b
r^2\ee^{-2r/\aB}\dd r\\
&=\boxed{1-\ee^{-2b/\aB}
\left(1+\frac{2b}{\aB}+
\frac{2b^2}{\aB^2}\right)}.
\end{align}
Let $\epsilon=2b/\aB\ll1$.
Expanding
$1-\ee^{-\epsilon}(1+\epsilon+\epsilon^2/2)
=\epsilon^3/6+O(\epsilon^4)$ gives
\begin{equation}
\boxed{P\simeq\frac43(b/\aB)^3}.
\end{equation}
Approximating $\psi$ by its value at the
origin in the nuclear volume gives
$(4\pi b^3/3)|\psi(0)|^2=
(4/3)(b/\aB)^3$, the same result.
With $b=10^{-15}$ m and
$\aB=0.5\times10^{-10}$ m,
\begin{equation}
\boxed{P\simeq1.07\times10^{-14}}.
\end{equation}
This uses a point-Coulomb wavefunction
inside the nuclear volume, as requested;
finite nuclear structure modifies it.

\problem{4.55}{Circular hydrogenic states and radial uncertainty}
For $\ell=n-1$, the radial polynomial
is constant, and
$R_{n,n-1}=N_n r^{n-1}\ee^{-r/(n\aB)}$.
Set $\beta=2/(n\aB)$.
Normalization requires
$1=|N_n|^2\int_0^\infty r^{2n}
\ee^{-\beta r}\dd r=
|N_n|^2(2n)!/\beta^{2n+1}$, so
\begin{equation}
\boxed{N_n=\frac{\beta^{n+1/2}}{\sqrt{(2n)!}}}.
\end{equation}
The radial moments follow from successive
gamma-function integrals:
\begin{align}
\boxed{\avg r}&=\frac{2n+1}{\beta}
=\frac{n\aB}{2}(2n+1),\\
\boxed{\avg{r^2}}&=\frac{(2n+1)(2n+2)}{\beta^2}.
\end{align}
Thus
$\sigma_r^2=(2n+1)/\beta^2$,
and
\begin{equation}
\boxed{\sigma_r=\frac{\avg r}{\sqrt{2n+1}}}.
\end{equation}
The fractional radial spread decreases
as $n^{-1/2}$ for large $n$,
suggesting a more localized radial
distribution in highly excited circular states.

\problem{4.56}{Two distinct hydrogen transitions with the same wavelength}
The Rydberg formula assigns
$\lambda^{-1}=R_H(n_f^{-2}-n_i^{-2})$.
Take the two different transitions
$(n_i,n_f)=(7,5)$ and $(35,7)$:
\begin{align}
\frac1{5^2}-\frac1{7^2}
&=\frac{49-25}{1225}=\frac{24}{1225},\\
\frac1{7^2}-\frac1{35^2}
&=\frac{25-1}{1225}=\frac{24}{1225}.
\end{align}
Both have $n_i>n_f$, hence
\begin{equation}
\boxed{\lambda_{7\to5}=\lambda_{35\to7}
=\frac{1225}{24R_H}}.
\end{equation}
This pair is not the example given
in the textbook.

\problem{4.57}{Uncertainty between $x^2$ and $L_z$}
Using $[L_z,x]=\ii\sh y$, one obtains
$[x^2,L_z]=-\ii\sh(xy+yx)
=-2\ii\sh xy$,
because $x$ and $y$ commute.
Robertson's inequality gives
\begin{equation}
\boxed{\sigma_{x^2}\,\sigma_{L_z}
\ge\sh|\avg{xy}|}.
\end{equation}
Every $\psi_{n\ell m}$ is an $L_z$
eigenstate, so $\sigma_{L_z}=0$.
Consequently
\begin{equation}
\boxed{\avg{xy}_{n\ell m}=0}.
\end{equation}
This agrees with azimuthal symmetry
of $|\psi_{n\ell m}|^2$.

\problem{4.58}{Probabilities for a complex spinor}
Take $\chi=A(1-2\ii,2)^T$.
Normalization gives $|A|^2(5+4)=1$;
choose $\boxed{A=1/3}$.
For $S_z$, the outcomes are
$\pm\sh/2$ with
\begin{equation}
\boxed{P_z(+)={5}/{9},\quad
P_z(-)={4}/{9},\quad
\avg{S_z}=\sh/18}.
\end{equation}
Since $a^*b=(1+2\ii)2/9=(2+4\ii)/9$,
$\avg{S_x}=\sh\,\mathrm{Re}(a^*b)=2\sh/9$,
$\avg{S_y}=\sh\,\mathrm{Im}(a^*b)=4\sh/9$.
The two-outcome relations
$P_i(\pm)=\frac12\pm\avg{S_i}/\sh$ give
\begin{align}
\boxed{P_x(+)=13/18,\ P_x(-)=5/18},\\
\boxed{P_y(+)=17/18,\ P_y(-)=1/18}.
\end{align}
Each pair sums to one and the
expectations follow from their weighted outcomes.

\problem{4.59}{Opposite spin correlations in the singlet}
The singlet is
$\ket{\mathrm s}=
(\ket{\uparrow\downarrow}
-\ket{\downarrow\uparrow})/\sqrt2$.
It satisfies
$\bra{\mathrm s}\sigma_i^{(1)}
\sigma_j^{(2)}\ket{\mathrm s}=-\delta_{ij}$:
the diagonal terms are $-1$ by rotational
invariance, and mixed terms vanish.
For unit vectors $\mathbf a,\mathbf b$,
\begin{align}
\avg{S_{\mathbf a}^{(1)}
S_{\mathbf b}^{(2)}}&=
\frac{\sh^2}{4}
\sum_{ij}a_ib_j
\avg{\sigma_i^{(1)}\sigma_j^{(2)}}\\
&=-\frac{\sh^2}{4}\,\mathbf a\cdot\mathbf b.
\end{align}
Writing $\mathbf a\cdot\mathbf b=\cos\theta$ gives
\begin{equation}
\boxed{\avg{S_{\mathbf a}^{(1)}
S_{\mathbf b}^{(2)}}=-\frac{\sh^2}{4}\cos\theta}.
\end{equation}

\problem{4.60}{Coupling spin one-half to arbitrary spin}
Let $s_2=j$ and set total $s=j\pm1/2$.
At fixed $m$, only two uncoupled states occur:
$\ket{+\frac12,j,m-\frac12}$ and
$\ket{-\frac12,j,m+\frac12}$.
Write
$\ket{sm}=A\ket{+\frac12,j,m-\frac12}
+B\ket{-\frac12,j,m+\frac12}$.
The matrix for $S^2$ on this span is
obtained from
$S^2=S_1^2+S_2^2+2S_{1z}S_{2z}+
S_{1+}S_{2-}+S_{1-}S_{2+}$.
Its diagonalization, followed by
$|A|^2+|B|^2=1$, yields
\begin{align}
\boxed{s=j+\tfrac12:}&\quad
A=\sqrt{\frac{j+m+1/2}{2j+1}},&
B=\sqrt{\frac{j-m+1/2}{2j+1}},\\
\boxed{s=j-\tfrac12:}&\quad
A=\sqrt{\frac{j-m+1/2}{2j+1}},&
B=-\sqrt{\frac{j+m+1/2}{2j+1}}.
\end{align}
This is the standard Condon--Shortley phase choice.
For $j=1/2$, these reproduce the triplet
and singlet coefficients; for $j=1$ and
$j=3/2$, substitution reproduces
the corresponding table entries.

\problem{4.61}{Spin-three-halves $S_x$ matrix and eigenvalues}
In the ordered $S_z$ basis
$m=3/2,1/2,-1/2,-3/2$,
\begin{equation}
\boxed{S_x=\frac\sh2
\begin{pmatrix}
0&\sqrt3&0&0\\
\sqrt3&0&2&0\\
0&2&0&\sqrt3\\
0&0&\sqrt3&0
\end{pmatrix}}.
\end{equation}
Writing $t=2\lambda/\sh$, the
characteristic polynomial is
\begin{equation}
\det(tI-2S_x/\sh)
=(t^2-1)(t^2-9).
\end{equation}
Therefore
\begin{equation}
\boxed{\lambda=\pm\frac{3\sh}{2},
\quad\pm\frac\sh2}.
\end{equation}
These eigenvalues agree with rotational
equivalence of $S_x$ and $S_z$.

\problem{4.62}{General spin-$s$ Cartesian matrices}
In the ordered $\ket{s,m}$ basis
$m=s,s-1,\ldots,-s$,
\begin{align}
\boxed{(S_z)_{m'm}}&=
\sh m\,\delta_{m'm},\\
\boxed{(S_+)_{m'm}}&=
\sh\sqrt{(s-m)(s+m+1)}\,
\delta_{m',m+1},\\
\boxed{(S_-)_{m'm}}&=
\sh\sqrt{(s+m)(s-m+1)}\,
\delta_{m',m-1}.
\end{align}
Hence
\begin{equation}
\boxed{S_x=\tfrac12(S_++S_-),\qquad
S_y=\tfrac1{2\ii}(S_+-S_-)}.
\end{equation}
They are $(2s+1)\times(2s+1)$ Hermitian
matrices with nonzero off-diagonal entries
only between adjacent $m$ values.
Their commutators obey
$[S_i,S_j]=\ii\sh\epsilon_{ijk}S_k$
and $S^2=\sh^2s(s+1)I$.
